% ==================================================================
% id:       func-009
% topic:    פונקציות
% subtopic: תחום הגדרה ותכונות של פונקציות
% solution: full
% ==================================================================
\textbf{עבור כל אחת מהפונקציות הבאות, קבעו האם הפונקציה מחזורית, ואם כן, חשבו את המחזור שלה.}
\begin{enumerate}
  \item{$f: \mathbb{R} \rightarrow \mathbb{R}, f(x)=x^2$}
\hint{הניחו שקיים מחזור $T>0$ ובדקו מה אומר התנאי $f(x+T)=f(x)$ בנקודה פשוטה, למשל $x=0$.}
\begin{solbox}
לא מחזורית
\end{solbox}
 %\item{$f: \mathbb{R} \rightarrow \mathbb{R}, f(x)=\sin(x)$}
 \item{$f: \mathbb{R} \rightarrow \mathbb{R}, f(x)=\sin(3x)$}
\begin{solbox}
 מחזורית, בעלת מחזור
 $T = \frac{2 \pi}{3}$.
\end{solbox}
 \item{$f: \mathbb{R} \rightarrow \mathbb{R}, f(x)=\cos(\frac{x}{4})$}
\begin{solbox}
 מחזורית, בעלת מחזור
 $T = 8 \pi$.
\end{solbox}
 \item{$f: \mathbb{R} \rightarrow \mathbb{R}, f(x)=\sin(3x) + \cos(\frac{x}{4})$}
\hint{מצאו את המחזור של כל מחובר בנפרד, ובדקו אם יש מספר חיובי שהוא כפולה שלמה של שניהם.}
\begin{solbox}
 מחזורית, בעלת מחזור
 $T = 8 \pi$.
\end{solbox}
 \item{$f: \mathbb{R} \rightarrow \mathbb{R}, f(x)=\lfloor x \rfloor$}

 פונקציה זו נקראת `''פונקציית הערך השלם'', והיא מתאימה לכל מספר ממשי את המספר השלם המקסימלי הקטן או שווה לו.
\hint{בדקו כיצד משתנים ערכי $\lfloor x\rfloor$ כאשר $x$ גדל: האם הפונקציה יכולה לחזור על ערכיה באופן מחזורי?}
\begin{solbox}
 לא מחזורית.
\end{solbox}
  \item{$f: \mathbb{R} \rightarrow \mathbb{R}, f(x)=\lfloor x \rfloor - x$}
\hint{בדקו מה קורה ל-$\lfloor x\rfloor$ ול-$x$ כאשר מוסיפים ל-$x$ מספר שלם.}
\begin{solbox}
 מחזורית, בעלת מחזור
 $T = 1$.
\end{solbox}
\end{enumerate}
